\documentclass[12pt]{article}

\usepackage{sbc-template}
\usepackage{graphicx,url}
\usepackage[utf8]{inputenc}
\usepackage[english]{babel}

\usepackage[color]{coqdoc}
\usepackage{mathpartir}
\usepackage{amssymb}
\usepackage{wasysym}
\usepackage{amsthm}
\newtheorem{definition}{Definition}
\newtheorem{lemma}{Lemma}
\newtheorem{theorem}{Theorem}
\newtheorem{proposition}{Proposition}
\newtheorem{corollary}{Corollary}

% Definitions are typeset in a breakable box with a colored left border.
\usepackage[skins,breakable]{tcolorbox}
\definecolor{DefColor}{RGB}{0,70,140}
\newtcolorbox{dfnbox}{%
  enhanced, breakable,
  colback=DefColor!4!white, colframe=DefColor!70!black,
  borderline west={3pt}{0pt}{DefColor},
  left=6pt, right=4pt%
}
\BeforeBeginEnvironment{definition}{\begin{dfnbox}}
\AfterEndEnvironment{definition}{\end{dfnbox}}
\usepackage{newunicodechar}
\newunicodechar{❬}{\ensuremath{\langle\!\langle}}
\newunicodechar{❭}{\ensuremath{\rangle\!\rangle}}
\newunicodechar{α}{\ensuremath{\alpha}}
\newunicodechar{∉}{\ensuremath{\notin}}
\newunicodechar{≠}{\ensuremath{\neq}}
\newunicodechar{∈}{\ensuremath{\in}}
\newunicodechar{∀}{\ensuremath{\forall}}
\newunicodechar{∖}{\ensuremath{\setminus}}
\newunicodechar{…}{\ensuremath{\dots}}

\newcommand{\alle}{(\forall_e)}
\newcommand{\alli}{(\forall_i)}
\newcommand{\toisa}{\rm{(\to_i)}}

\newcommand{\vswap}[3]{(\!\!(#1\ #2)\!\!) #3}
\newcommand{\swap}[3]{(#1\ #2) #3}
%\newcommand{\msub}[2]{\{#1 / #2\}}
\newcommand{\esub}[3]{[ #2 := #3 ] #1}
\newcommand{\metasub}[3]{\{ #2 := #3 \}#1}
\newcommand{\graft}[3]{\{ #2 /\!\!/ #3 \}#1}
\newcommand{\renv}[3]{\{ #2 \leadsto #3 \}#1}

\newcommand{\tto}{\twoheadrightarrow}
\newcommand{\ott}{\twoheadleftarrow}

\newcommand{\flavio}[1]{{\color{red}#1}}
\newcommand{\luis}[1]{{\color{blue}#1}}
     
\sloppy

\title{Formalizing \(\alpha\)-equivalence equations}
\author{Flávio L. C. de Moura\inst{1}, Luís A. M. Sales\inst{1}}

\address{Departamento de Ciência da Computação\\
  Universidade de Brasília -- Brasília, DF -- Brazil
  \email{flaviomoura@unb.br, sales.luis@aluno.unb.br}
}
 
\begin{document}

\maketitle

\begin{abstract}
  Syntax with binding constructors is pervasive in Computer Science and Mathematics, particularly in the formal specification and implementation of logics and programming languages. Characterizing terms that differ only by the names of their bound variables, i.e., terms in the same \(\alpha\)-equivalence class, is of great importance. A number of different definitions of \(\alpha\)-equivalence can be found in the current literature, employing various approaches, such as nameless representations (e.g., de Bruijn indices) or explicit-name approaches based on substitution/renaming or swapping. In this paper, we prove the equivalence between several definitions of \(\alpha\)-equivalence in the context of named variables using the Rocq proof assistant.  
\end{abstract}

\input{aeq_nexp.v}
\bibliographystyle{sbc}
\bibliography{referencesSep2026}
\end{document}
